\documentclass[10pt]{article}
\usepackage[margin=20mm]{geometry}
\usepackage{amsmath,amssymb,mathtools,fontspec,unicode-math,xurl,hyperref}
\IfFontExistsTF{DejaVu Serif}{\setmainfont{DejaVu Serif}}{\setmainfont{Latin Modern Roman}}
\setmathfont{latinmodern-math.otf}
\setlength{\parindent}{0pt}\setlength{\parskip}{6pt}\sloppy
\title{Technical review of
Unknotting nonorientable surfaces}
\author{AI technical verification}
\date{}
\begin{document}
\maketitle
Abstract: The assessment of this accepted author manuscript is Major. An exact rank-two example satisfies every hypothesis of the central lagrangian-complement lemma but disproves the basis construction used in its proof. The failure affects the stated route from the plus and minus surgery obstructions to the main classification results. A complete replacement proof of the algebraic lemma is provided, so the finding does not assert that the lemma or the main classification theorems are false. The Euler-characteristic inference used to establish simple connectedness in Corollary D is a second Major proof defect, with a meridian-monodromy repair. Six separate Minor statements require restriction or correction. All 64 physical source pages were read, and the recorded evidence pages were inspected as images. The source version, finite computations, external dependencies and limits of verification are distinguished throughout. Scientific adjudication remains pending.\par
\section*{Source identity and review scope}
The assigned paper is M25-0265, Unknotting nonorientable surfaces, by Anthony Conway, Patrick Orson and Mark Powell, DOI `10.4171/jems/1638`. The reviewed file is the University of Glasgow repository’s accepted author version, item 340906: 63 article pages followed by a repository notice, 918,302 bytes. Physical and printed article page numbers agree. Its SHA256 is\par
> de8f8db50ba80e75796549812f9ac29295b8e39226bff89fa9d9beba400663ad.\par
The notice records deposit on 6 November 2024. The repository labels the file Accepted Version. The EMS record gives submission on 13 July 2023, acceptance on 20 October 2024 and online publication on 20 May 2025. The associated arXiv record has version 2 dated 23 November 2024. That metadata does not establish that the repository PDF, the arXiv PDF and the publisher PDF are textually identical. The file’s 2026 container timestamps are not used as a scientific revision date. The frozen title, authors, DOI and year are retained.\par
The principal statements concern locally flat nonorientable \(\mathbb Z_2\)-surfaces: unknottedness in the specified genus and normal Euler-number range, classification relative to a determinant-one boundary knot, stabilization by one orientable tube, branched covers, and realization of a prescribed Guillou–Marin isometry. The review checked the precise range and boundary hypotheses, the homology and intersection forms, the spin gluing, the modified surgery obstruction, the plus/minus algebra, the cancellation step, and the finite discriminant computations.\par
The assessment comprises two Major and six Minor findings. It includes the essential construction in Lemma 9.12 and the simple-connectedness inference used in Corollary D. Their repairability does not remove their importance to the submitted derivations. The other numbered findings have the specific local scope and counterevidence given below. No counterexample to Theorems A–E is claimed. Full reading is not a claim to have reproved the entire external surgery and four-manifold literature.\par
\section*{The central algebraic construction}
Problem 1 (Major: the dual vectors in Lemma 9.12 need not complete a basis).  On p. 50, after completing a basis of \(V\) to one of \(P\), the proof chooses \(u_i\) by \(\widehat\lambda(u_i)=x_i^*\) and asserts that the \(x_i,u_i\) generate \(P\). This assertion is false under all the stated hypotheses. The error occurs when an integral specialization of a matrix identity is treated as an identity over \(\mathbb Z[\mathbb Z_2]\).\par
\subsection*{An exact example satisfying every hypothesis}
Put \(\Lambda=\mathbb Z[T]/(T^2-1)\) with the involution used in the paper, and \(P=\Lambda x\oplus\Lambda v\). Define the quadratic representative and its symmetrization by \[Q=\begin{pmatrix}2(1-T)&4-3T\\0&3\end{pmatrix},
 \qquad
 A=Q+Q^t=
 \begin{pmatrix}4(1-T)&4-3T\\4-3T&6\end{pmatrix}.\] In \(\Lambda\), \[\det A=24(1-T)-(4-3T)^2=-1.\] Thus the hermitian form is nonsingular and even. The module \(V=\Lambda x\) is a free half-rank direct summand. In the natural bases of the eigenspaces, Construction 9.3 gives precisely the specialized matrices \[A_+=\begin{pmatrix}0&1\\1&6\end{pmatrix},
 \qquad
 A_-=\begin{pmatrix}8&7\\7&6\end{pmatrix}.\] The plus submodule \(\mathbb Zx\) is a lagrangian. In the minus form, \(\mathbb Z(v-x)\) is a direct complement to \(\mathbb Zx\) and is isotropic because \(8-14+6=0\). Its pairing with \(x\) is \(-1\), so it is a lagrangian. These verify every hypothesis of Lemma 9.12, including the integral direct-summand conditions.\par
Use the allowed basis completion \(x,v\). The uniquely prescribed dual vector is \[u=A^{-1}x^*=-6x+(4-3T)v.\] It satisfies \(\lambda(x,u)=1\) and \(\lambda(v,u)=0\), as required. But \[\det[x,u]=4-3T,\qquad (4-3T)|_{T=-1}=7,\] so \([x,u]\) is not invertible over \(\Lambda\). More concretely, \[P/(V+\Lambda u)\cong \Lambda/(4-3T).\] On the additive basis \(1,T\), multiplication by \(4-3T\) has matrix \(\left(\begin{smallmatrix}4&-3\\-3&4\end{smallmatrix}\right)\), whose Smith invariants are \(1,7\). The quotient is therefore \(\mathbb Z/7\), and the asserted generation fails.\par
This example also has the relevant restricted quadratic form: \(Q|_V=2(1-T)=(1-T)X_1\) in the paper’s notation. Moreover, \[\ell=v-(2+T)x,\qquad Q(\ell)=0,\qquad
 \lambda(x,\ell)=T,\qquad \det[x,\ell]=1.\] Thus \(\Lambda\ell\) is an actual lagrangian complement, and the example can occur as the algebra of an elementary quasi-formation. It disproves the proposed construction, not the conclusion of the lemma.\par
The exact calculation was independently reproduced using self-written polynomial arithmetic modulo \(T^2-1\) and integer Smith normal form. A simpler companion example with diagonal \(2(1-T),2\) and cross entry \(2-T\) gives quotient \(\mathbb Z/3\). The first check script contained an incorrect expected cross-pairing for that companion; the actual value is \(T\). Its failed exit and original script are preserved. The corrected script exited successfully and checked both examples. No author code was run.\par
\subsection*{The invalid identity and its effect}
With the notation \(u=Ax+Bv\) of p. 50, the valid specialized identity is \[B(+1)L=I.\] The proof writes \(BL=I\) over \(\Lambda\). Specialization replaces \(b_{ij}\) by \(b_{ij}(+1)\); it cannot justify invertibility of \(B\) over the group ring. The example above has \(B=4-3T\) and \(B(+1)=1\).\par
There is a further manifestation of the same construction problem. On p. 51 the proof changes \(U_+\) by adding vectors from \(\mathcal E_+(V)\) and expressly notes that \(U_+=\mathcal E_+(U_1)\) may no longer hold. On p. 52 it invokes that identity again to place \(c_i\) in \(U_1\) and infer a new basis. Neither the containment nor invertibility of a general modification follows from the preceding text. These observations are grouped with Problem 1 rather than counted as independent failures of the classification theorem.\par
Lemma 9.12 is used by Proposition 9.14, then Proposition 9.15 and the obstruction-vanishing argument in Theorem 10.1. The failure is therefore in the core proof chain, not an unused generality. The existence of the repair below limits the conclusion: the printed route is invalid, while this algebraic obstruction to the route can be removed.\par
\subsection*{A complete replacement proof of Lemma 9.12}
Choose any free complement \(U_0\) to \(V\). Freeness is available from Lemma 3.9. Relative to \(\mathcal E_-(V)\oplus\mathcal E_-(U_0)\), the given complementary minus lagrangian is the graph of an integral matrix. Lift that matrix through \(\mathbb Z\hookrightarrow\Lambda\) and replace \(U_0\) by the corresponding graph \(U\). Then \(P=V\oplus U\) and \(\mathcal E_-(U)\) is the given lagrangian.\par
In bases of \(V,U\), write \[\lambda=\begin{pmatrix}A&Y\\Y^t&D\end{pmatrix}.\] We have \(A(+1)=0\) and \(D(-1)=0\). Nonsingularity of the specialized forms and the block determinant formula give \(\det Y(+1),\det Y(-1)\in\{1,-1\}\). Identify \[\Lambda\cong\{(a_+,a_-)\in\mathbb Z^2:a_+\equiv a_-\pmod2\}.\] The integral matrices \(Y(+1),Y(-1)\) agree modulo 2. Their inverses therefore agree modulo 2 as well and lift to an inverse over \(\Lambda\). Change the basis of \(U\) using \(Y^{-1}\), so the mixed block is \(I\). The vanishing specializations imply \[A=(1-T)A_0,\qquad D=(1+T)D_0\] for integral symmetric matrices \(A_0,D_0\). Evenness of the ambient form makes every diagonal entry of \(D_0\) even. Hence there is an integral matrix \(C\) with \(C+C^t=-D_0\): take \(C_{ii}=-(D_0)_{ii}/2\), the required negative entries above the diagonal, and zero below.\par
Use the graph map \(u_j\mapsto u_j+\sum_i v_i(1+T)C_{ij}\). This remains a free complement. Its Gram block is \[D+(1+T)(C+C^t)+(1+T)^2C^tAC=0\] because \((1-T)(1+T)=0\). The mixed block remains \(I\) for the same reason. The new complement is therefore its own annihilator. It is a quadratic lagrangian as well, since symmetrization is injective over this particular torsion-free ring with trivial involution. This proves the stated lemma without the invalid dual-basis construction or the inconsistent plus-basis changes.\par
\section*{Further statements and proof steps requiring correction}
Problem 2 (Minor: symmetrization is not injective over every stated ring).  Lemma 3.7 and Corollary 3.8, p. 11, are stated for a general ring with involution. Take \(R=\mathbb F_2\) with trivial involution, \(\varepsilon=1\) and the free module \(P=R\). Then \(Q_+(P)=\mathbb F_2\), while symmetrization is multiplication by 2 and is zero. The nonzero class represented by \([1]\) lies in its kernel.\par
The zero bilinear form has the two different refinements \(\mu_0(x)=0\) and \(\mu_1(x)=x^2\). They satisfy all three conditions of Definition 3.3. With \(L=P\), the bilinear form vanishes on \(L\) but the nonzero quadratic form does not. Thus both the uniqueness and the claimed isotropic-to-sublagrangian implication fail in this generality. In the proof, a skew-hermitian matrix need not have zero diagonal under the stated hypotheses.\par
The actual applications here use \(\mathbb Z\) and \(\mathbb Z[\mathbb Z_2]\) with trivial involution and no additive 2-torsion. For these rings the diagonal obstruction vanishes, the triangular argument works, and the projective reduction works. Restrict the lemma and corollary to that valid setting, or give the appropriate general hypothesis on symmetrization. This is a local overstatement; it does not invalidate the intended applications.\par
Problem 3 (Minor: stable freeness alone does not imply the surjectivity in Remark 3.6).  On p. 11, let \(R=\mathbb Z\), \(P=L=\mathbb Z\), \(\psi=0\), and let \(j\) be the identity. All modules are free, \(L\) is a direct summand and \(L=L^\perp\) under Definition 3.5. Nevertheless \(j^*(1+\mathord\top)\psi:P\to L^*\) is the zero map and is not surjective.\par
The needed extra condition is ambient nonsingularity, not stable freeness. With that condition the map to \(P^*\) is an isomorphism and restriction \(P^*\to L^*\) is onto for a direct summand. The counterexample is excluded from the nonsingular quasi-formations used later. Correct the remark and the scope of the lagrangian convention; no additional failure of the main surgery argument is inferred.\par
Problem 4 (Minor: the excision assertion in the proof of Proposition 4.2 has excessive degree range).  On p. 13 the proof identifies the relative homology of the full disc/sphere bundle pair with that over the punctured surface for all degrees. Let \(F=\mathbb{RP}^2\) and \(\xi=TF\). This is a rank-two bundle with \(w_1(\xi)=w_1(F)\), so its total disc bundle is oriented, as required. For its disc/sphere pair \((N,Y)\), the Thom isomorphism with the orientation local system gives \[H_4(N,Y;\mathbb Z)=H_2(F;\mathbb Z_{w_1})=\mathbb Z.\] For the punctured surface \(\dot F\), \[H_4(\dot N,\dot Y;\mathbb Z)
   =H_2(\dot F;\mathbb Z_{w_1})=0.\] The stated all-degree isomorphism is false.\par
Only degrees 1 and 2 are used at this point. They are respectively zero and \(\mathbb Z/2\) in both cases. Applying the Thom isomorphism in those degrees proves exactly what the subsequent homology argument needs. The counterexample refutes the intermediate universal excision assertion, not Proposition 4.2’s required conclusion.\par
Problem 5 (Minor: global degree is misstated in Lemma 4.5).  On p. 15, the proof calls the branched double-cover maps degree one. A regular value off the branch locus has two orientation-preserving preimages. The global topological degree is therefore two. Each unbranched sheet has local degree one, and these two notions must be distinguished.\par
The correct fibre relation gives the stated normal Euler-number identity \(e=2e'\). The global degree-one explanation should be replaced by the local sheet calculation and the squaring map on the normal fibre. This correction preserves the lemma and the later signature formula.\par
Problem 6 (Major: the mod-2 Euler-characteristic argument for Corollary D).  In the proof on pp. 3–4, if the hypothetical residual fundamental group were \(\mathbb Z_2\), then \(b_1(X;\mathbb F_2)=b_3(X;\mathbb F_2)=1\) by abelianization and Poincare duality. Consequently \[\chi(X)=1-1+b_2(X;\mathbb F_2)-1+1=b_2(X;\mathbb F_2),\] not \(2+b_2(X;\mathbb F_2)\). The claimed contradiction from the lower bound on the mod-2 second Betti number does not follow.\par
For a concrete counterexample to that auxiliary Euler inference, the free diagonal antipodal quotient \[M=(S^2\times S^2)/((x,y)\sim(-x,-y))\] is closed and oriented, has fundamental group \(\mathbb Z_2\), and satisfies \(\chi(M)=2\) and \(b_2(M;\mathbb F_2)=2\). It is not asserted to be the branched cover of a surface meeting every hypothesis of Corollary D.\par
The invalid inequality is the printed justification for simple connectedness of the orbit space. This conclusion is then needed to identify that space with \(S^4\) by Freedman’s theorem and apply Theorem A. The defect is therefore Major for the derivation of this stated involution application. The auxiliary quotient above does not refute the corollary under all of its branched-cover hypotheses.\par
A valid repair uses local monodromy: a meridian has nontrivial image in the deck group of the branched double cover, hence generates the exterior group \(\mathbb Z_2\). The van Kampen quotient kills that generator and establishes simple connectedness without the erroneous Betti-number inequality.\par
Problem 7 (Minor: the exterior quadratic function cannot be extended linearly).  Definition 6.14 on p. 30 says to extend the half-twist prescription linearly from a basis of \(H_1(F;\mathbb F_2)\) to \(\mathbb Z_4\). Any additive map \(\mathbb F_2\to\mathbb Z_4\) sends its generator to 0 or 2. A one-sided generator has self-intersection 1 and its Guillou–Marin value is odd. For example the standard punctured projective plane in \(D^4\) has such a generator. Thus no additive extension can have the required value.\par
Use the quadratic extension \(q(x+y)=q(x)+q(y)+2Q_F(x,y)\), with the coefficient conventions of Section 6. Proposition 6.15 supplies the correct geometric identification with the Guillou–Marin form and hence repairs the definition. The disc framing used later is expressly required to extend over the disc on p. 31; this review does not claim that requirement is absent.\par
Problem 8 (Minor: the genus-three appendix changes its scale).  On pp. 57 and 60 the statement of Proposition A.6 displays \[\theta=2X_3=
 \begin{pmatrix}4&4&4\\0&2&2\\0&0&2\end{pmatrix}\] but asks about \(\operatorname{bAut}(V,2\theta)\) and sets \(\psi=2\theta\). These expressions refer to \(4X_3\). The proof instead uses the discriminant form of \(2X_3\) from Lemmas A.4–A.5.\par
The distinction is detected without numerical approximation: \(\det(2X_3+2X_3^t)=32\), whereas \(\det(4X_3+4X_3^t)=256\). Their discriminant groups have orders 32 and 256 respectively; in Smith form they are \(\mathbb Z_8\oplus\mathbb Z_2^2\) and \(\mathbb Z_{16}\oplus\mathbb Z_4^2\). Thus the displayed group and quadratic values used in the proof cannot be those of its literally defined \(\psi\). A further exact check shows that the literally doubled assertion is false, not merely unsupported. In the coordinates of Lemma A.4 the boundary of \(4X_3\) is \(\mathbb Z_{16}\oplus\mathbb Z_4^2\) with \[q(a,b,c)=\frac{11a^2-16ab+8b^2+8bc+8c^2}{32}\pmod{\mathbb Z}.\] Consider the well-defined automorphism \[U(a,b,c)=(a+8c,\ b+c,\ 3c).\] It has square the identity on this finite group, and \[q(U(a,b,c))-q(a,b,c)=5ac-3bc+21c^2\in\mathbb Z.\] It is therefore a quadratic boundary isometry.\par
Here is a direct proof that it does not lift. Realize the original integral lattice as the even-coordinate-sum sublattice \(L\subset\mathbb Z^3\), with basis \(2e_1,e_1+e_2,e_1+e_3\) and the Euclidean pairing multiplied by 4. The vectors of Euclidean squared length 4 in \(L\) are precisely \(\{\pm2e_i\}\), so every integral isometry is a signed coordinate permutation. With the unimodular coordinate change \(D=BC\) printed in Lemma A.4, the first two boundary generators have dual representatives \[w_1=(-1,-1,-3)/8,\qquad w_2=(1,0,1)/4.\] A signed permutation fixing the class of \(w_1\) must fix \(w_1\) itself: each coordinate of its difference from \(w_1\) has absolute value less than 1, while a vector of \(L\) has integral coordinates. The only possibilities are the identity and interchange of the first two coordinates. The latter does not fix \(w_2\) modulo \(L\). Thus only the identity fixes both boundary generators. But \(U\) fixes both and is nonidentity. It cannot be induced by an integral isometry. Its double coset is consequently different from the identity double coset in the stated boundary-automorphism set.\par
A separate exact enumeration finds 384 boundary isometries and 48 distinct integral images; it checks all 256 group elements for the displayed \(U\). The analytic argument above supplies the decisive non-lift proof without relying on the counts. For the intended genus-three application in Proposition 7.44, one needs \(\theta=X_3\) and the boundary of \(2\theta\). Correcting the displayed scale restores the object treated by Lemmas A.4–A.5. This remains a Minor finding because the explicit counterexample addresses the incorrectly doubled auxiliary statement, not the form required by the main classification.\par
\section*{Further checks and limits of the assessment}
\subsection*{Ranges, topology and geometric compatibility}
The ranges of A and B are the low-genus or nonextremal normal-Euler-number cases; the determinant-one condition belongs to B, while the one-tube stabilization in C has its separately stated hypotheses. Theorem E retains the prescribed Guillou–Marin isometry, and its gluing proof must preserve that choice. These distinctions were checked through Sections 10–11.\par
The displayed Massey range in Theorem 2.4 on p. 9 is malformed. The intended set, consistent with pp. 1 and 10, is \(\{-2h+4j:0\le j\le h\}\). Connected sums of \(a\) positive and \(b\) negative standard projective planes realize the corresponding endpoints and steps with the paper’s sign convention. Example 2.3’s mention of a single 0-handle and a single 1-handle does not explicitly say these are the only handles. To deduce \(\pi_1=\mathbb Z_2\) from a handle description, the square-meridian 2-handle relation should also be supplied. This is recorded as a clarification, not as a false claim that no 2-handles occur.\par
The boundary and branched-cover homology calculations in Sections 4–5 were checked with the twisted coefficients and Euler characteristic retained. The index-two even lattice calculation and the special \(\sigma=0\) hyperbolic case are consistent. In the definition of \(X_n\) for negative \(n\), the rank is \(|n|\), as the actual matrices and applications require.\par
The Section 6 construction of nice sections, the use of immersed discs in the simply connected cover, the odd-framing comparison, and the spin structures extending over the exteriors were followed. The line-preserving derivative in Lemma 6.20(2) need not already be the identity: a diagonal linear map is an immediate local example. A positive normalization after fixing the sign provides the required adjustment. No all-hypothesis counterexample to the extension lemma was established.\par
\subsection*{Surgery algebra and finite computations}
The plus/minus normalization of Construction 9.3 was checked directly: evaluating a matrix at \(T=\pm1\) is consistent with division by \(2(1\pm T)\) on the eigenspace pairings. The nondegenerate Wall form and its minus specialization produce the factor \(2\theta\) used in the integral obstruction. Zero-stable equivalence, formations, quasi-formations and the boundary automorphism double-coset set were kept distinct.\par
For Proposition 7.44, the discriminant is 2-primary, so the odd-prime generator counts are zero. The rank inequalities in the indefinite case and after stabilization by \(2H_+(\mathbb Z)\) were checked. The small definite cases require the separate appendix; one cannot use the indefinite theorem there. Independent exact calculations for ranks \(1\) through \(16\) recover the Smith pattern \[\mathop{\mathrm{coker}}(2X_n+2X_n^t)\cong
 \begin{cases}
  \mathbb Z_2^{\,n-1}\oplus\mathbb Z_8,&n\text{ odd},\\
  \mathbb Z_2^{\,n-2}\oplus\mathbb Z_4\oplus\mathbb Z_4,&n\text{ even}.
 \end{cases}\] Those finite checks support the algebra and were not presented as a proof for every rank. The index-two parity-lattice description gives the general structural explanation. Enumeration of the rank-two quadratic boundary form gives eight isometries, while the \(4H_+\) boundary form has four. These calculations were made with exact integers and rational values, not floating point.\par
Remark 7.45 on p. 44 says that no nontrivial extremal examples are known, whereas the introduction cites the newer Pen24 examples in genera six and seven. This is a stale cross-reference in the reviewed version, not a further refutation of a theorem.\par
\subsection*{External dependencies and unavailable material}
The proof uses deep external results: realization of quadratic enhancements, topological spin bordism, Kreck’s modified surgery, the good-group topological \(s\)-cobordism theorem, integral lattice classification, Nikulin’s boundary-isometry criterion, and the stated cancellation theorem. Their precise uses and hypotheses were traced in the paper; a complete independent reconstruction of those libraries was not performed. In particular, reading the bordism group identification and the cancellation citation is not equivalent to redoing all spectral-sequence differentials or the general cancellation proof.\par
Bounded searches of the title, DOI, correction and supplementary terms, together with the EMS, Glasgow and arXiv primary records, did not identify an applicable separate correction or public code package. This is not a proof of absence. The appendix names Sage computations and says matrices are available on request. No request was sent and no author code was executed. Lack of public code is not treated as a mathematical defect.\par
\section*{Verification record and recommendation}
The initial primary reviewer read all 64 physical text pages and actually viewed the source images listed in the saved coverage record. Three native AI auxiliaries provided bounded checks within this original review workflow; their individual scope, runtime self-description, inspected pages and limitations are retained. Their initial work did not constitute independent acceptance. The current focused independent acceptance and its read-only auxiliary contribution are separately attributed in the verification records. No human review is claimed. The reported UI verification belongs to the external controller; the runtime’s own identity statement is separately recorded and is not an independent backend attestation.\par
Rendering alone is not counted as viewing. Every current report page is subject to actual visual inspection before delivery is marked complete.\par
The recommended mathematical action is to replace the proof of Lemma 9.12, including its downstream citations to the repaired argument, to replace the Corollary D Euler inference by the meridian argument, and to correct the local generality, degree, coefficient and scale statements identified above. The replacement proof here settles that lemma under its actual hypotheses; it does not constitute a new independent verification of every geometric dependency of A–E. The assessment is two Major and six Minor findings, with `scientific\_decision=pending` and no human reviews recorded.\par
9 A. Conway, P. Orson and M. Powell, Unknotting nonorientable surfaces, accepted author manuscript, Glasgow repository item 340906, exact source reviewed here, https://eprints.gla.ac.uk/340906/2/340906.pdf. EMS Press, article record, https://ems.press/journals/jems/articles/14298811. University of Glasgow, accepted-version metadata and repository notice, https://eprints.gla.ac.uk/340906/. A. Conway, P. Orson and M. Powell, arXiv record 2306.12305, version history, https://arxiv.org/abs/2306.12305. The record was used for version metadata, without assuming full-text equivalence.\par
\end{document}
